% ============================================================
% f2012_guangdong_arts_q07_views.tex
% 2012 广东卷（文科）第 7 题 三视图（TikZ 矢量重绘）
% 正视图、侧视图：上半为半径 3 的半圆，下半为底 6、腰 5（高 4）的倒三角
% 俯视图：直径 6 的圆；还原为半球（r=3）与圆锥（r=3, h=4）的组合体
% ============================================================
\begin{tikzpicture}[
  x=1cm, y=1cm, line width=0.6pt, font=\footnotesize, >=Stealth
]

  % ---------------- 正视图 ----------------
  \begin{scope}
    % 半圆（上）
    \draw (-1.5,0) arc[start angle=180, end angle=0, radius=1.5];
    \draw (-1.5,0) -- (1.5,0);
    % 三角形（下，顶点朝下）
    \draw (-1.5,0) -- (0,-2) -- (1.5,0);
    % 上方尺寸 6
    \draw[<->] (-1.5,1.85) -- (1.5,1.85);
    \draw (-1.5,1.78) -- (-1.5,1.92);
    \draw (1.5,1.78) -- (1.5,1.92);
    \node[above] at (0,1.90) {$6$};
    % 右侧尺寸 3
    \draw[<->] (1.78,0) -- (1.78,1.5);
    \draw (1.72,0) -- (1.86,0);
    \draw (1.72,1.5) -- (1.86,1.5);
    \node[right] at (1.80,0.75) {$3$};
    % 两腰尺寸 5
    \node[rotate=-53.13] at (-1.12,-1.24) {$5$};
    \node[rotate=53.13] at (1.12,-1.24) {$5$};
    % 视图名
    \node at (0,-2.55) {正视图};
  \end{scope}

  % ---------------- 侧视图 ----------------
  \begin{scope}[xshift=4.9cm]
    \draw (-1.5,0) arc[start angle=180, end angle=0, radius=1.5];
    \draw (-1.5,0) -- (1.5,0);
    \draw (-1.5,0) -- (0,-2) -- (1.5,0);
    \draw[<->] (-1.5,1.85) -- (1.5,1.85);
    \draw (-1.5,1.78) -- (-1.5,1.92);
    \draw (1.5,1.78) -- (1.5,1.92);
    \node[above] at (0,1.90) {$6$};
    \draw[<->] (1.78,0) -- (1.78,1.5);
    \draw (1.72,0) -- (1.86,0);
    \draw (1.72,1.5) -- (1.86,1.5);
    \node[right] at (1.80,0.75) {$3$};
    \node[rotate=-53.13] at (-1.12,-1.24) {$5$};
    \node[rotate=53.13] at (1.12,-1.24) {$5$};
    \node at (0,-2.55) {侧视图};
  \end{scope}

  % ---------------- 俯视图 ----------------
  \begin{scope}[yshift=-5.8cm]
    \draw (0,0) circle (1.5);
    \node at (0,-1.95) {俯视图};
  \end{scope}

\end{tikzpicture}
